\documentclass[11pt,a4paper]{article}

% ── Packages ──────────────────────────────────────────────────────────────────
\usepackage[T1]{fontenc}
\usepackage[utf8]{inputenc}
\usepackage{geometry}
\usepackage{hyperref}
\usepackage{xcolor}
\usepackage{listings}
\usepackage{longtable}
\usepackage{booktabs}
\usepackage{array}
\usepackage{enumitem}
\usepackage{titlesec}
\usepackage{fancyhdr}
\usepackage{parskip}
\usepackage{microtype}
\usepackage{lmodern}

\geometry{margin=2.4cm}

% ── Colors ────────────────────────────────────────────────────────────────────
\definecolor{codebg}{HTML}{F5F5F5}
\definecolor{codeborder}{HTML}{CCCCCC}
\definecolor{kw}{HTML}{0000CC}
\definecolor{str}{HTML}{008800}
\definecolor{comment}{HTML}{888888}
\definecolor{typecolor}{HTML}{7B0000}
\definecolor{linkcolor}{HTML}{0066CC}
\definecolor{tableheadbg}{HTML}{E8EEF4}
\definecolor{sectionrule}{HTML}{3A6EA5}

% ── Hyperref setup ─────────────────────────────────────────────────────────────
\hypersetup{
  colorlinks=true,
  linkcolor=linkcolor,
  urlcolor=linkcolor,
  citecolor=linkcolor,
  pdftitle={UnityIRRuntime --- Feature \& API Reference},
  pdfauthor={ObjectIR.Unity.Runtime},
}

% ── Code listing style ────────────────────────────────────────────────────────
\lstdefinestyle{cs}{
  language=[Sharp]C,
  backgroundcolor=\color{codebg},
  frame=single,
  rulecolor=\color{codeborder},
  basicstyle=\ttfamily\small,
  keywordstyle=\bfseries\color{kw},
  stringstyle=\color{str},
  commentstyle=\itshape\color{comment},
  breaklines=true,
  showstringspaces=false,
  tabsize=4,
  numbers=none,
}

\lstdefinestyle{inline}{
  basicstyle=\ttfamily\small,
  keywordstyle=\bfseries\color{kw},
  breaklines=true,
}

\newcommand{\code}[1]{\texttt{\small #1}}
\newcommand{\typename}[1]{\textcolor{typecolor}{\texttt{\small #1}}}
\newcommand{\kword}[1]{\textbf{\texttt{\small #1}}}

% ── Headers / footers ─────────────────────────────────────────────────────────
\pagestyle{fancy}
\fancyhf{}
\renewcommand{\headrulewidth}{0.4pt}
\fancyhead[L]{\small\textit{UnityIRRuntime --- Feature \& API Reference}}
\fancyhead[R]{\small\textit{ObjectIR.Unity.Runtime}}
\fancyfoot[C]{\thepage}

% ── Section rule decoration ───────────────────────────────────────────────────
\titleformat{\section}{\large\bfseries\color{sectionrule}}{}{0em}{}[{\color{sectionrule}\titlerule[0.8pt]}]
\titleformat{\subsection}{\normalsize\bfseries}{}{0em}{}
\titleformat{\subsubsection}{\normalsize\itshape\bfseries}{}{0em}{}

% ── Document ──────────────────────────────────────────────────────────────────
\begin{document}

\begin{titlepage}
  \centering
  \vspace*{3cm}
  {\Huge\bfseries UnityIRRuntime}\\[0.6em]
  {\Large Feature \& API Reference}\\[2em]
  {\large\texttt{ObjectIR.Unity.Runtime}}\\[0.4em]
  {\normalsize namespace \texttt{ObjectIR.Unity.Runtime}}\\[3em]
  \rule{0.5\linewidth}{0.4pt}\\[1em]
  {\normalsize
    Source files: \texttt{UnityIRRuntime.cs}, \texttt{UnityIRRuntime.Types.cs}\\[0.3em]
    Assembly definition: \texttt{Gamerboy030607.Objectirunity}\\[0.3em]
    Unity package: \texttt{UnityIR/Runtime}
  }
  \vfill
  {\small\today}
\end{titlepage}

\tableofcontents
\newpage

% ═══════════════════════════════════════════════════════════════════════════════
\section{Overview}
% ═══════════════════════════════════════════════════════════════════════════════

\textbf{UnityIRRuntime} is a Unity-aware, standalone ObjectIR virtual machine. It
embeds the full IR execution engine together with a comprehensive set of
\texttt{UnityEngine} API bindings so that ObjectIR scripts can drive Unity
\emph{without} any hand-written glue code for common APIs.

Key design goals:

\begin{itemize}[leftmargin=1.5em]
  \item \textbf{Zero external assembly dependency at runtime.} All IR type definitions
    (\typename{ModuleDto}, \typename{TypeDto}, \typename{MethodDto}, etc.) are
    embedded directly in \texttt{UnityIRRuntime.Types.cs}.
  \item \textbf{Multiple input formats.} A module can be loaded from JSON, TextIR
    source, a pre-parsed AST, or a compiled FOB/IR v3 binary.
  \item \textbf{Reflection fallback.} Any Unity type that is not covered by a manual
    binding is handled automatically via CLR reflection, so even obscure APIs
    work out of the box.
  \item \textbf{Hot-path manual bindings.} Frequently called APIs
    (\typename{Debug}, \typename{Vector3}, \typename{Transform}, \typename{Input},
    \typename{Mathf}, \ldots) have hand-written, allocation-light bindings that
    bypass reflection entirely.
  \item \textbf{Struct boxing/unboxing.} Unity value types
    (\typename{Vector2}/\typename{3}, \typename{Quaternion}, \typename{Color}) are
    transparently wrapped in \typename{ManagedObject} instances so the IR heap
    can hold them without loss of precision.
\end{itemize}

\subsection{Quick-start}

\begin{lstlisting}[style=cs]
// 1. Create the runtime and tie it to a GameObject.
var rt = new UnityIRRuntime(gameObject);

// 2. Load a module (TextIR, JSON, or binary).
rt.LoadModule(scriptText);          // auto-detects format
rt.LoadModule(bytes);               // raw FOB/IR binary
rt.LoadFobFile("path/to/mod.fob");  // from disk

// 3. Invoke methods by qualified name.
rt.CallMethod("Program.OnStart");
rt.CallMethod("Program.OnUpdate", Time.deltaTime);
float hp = rt.CallMethod<float>("Player.GetHealth");

// 4. Check whether an optional lifecycle method exists before calling.
if (rt.HasMethod("Program.OnFixedUpdate"))
    rt.CallMethod("Program.OnFixedUpdate");
\end{lstlisting}

% ═══════════════════════════════════════════════════════════════════════════════
\section{Class Declaration}
% ═══════════════════════════════════════════════════════════════════════════════

\begin{lstlisting}[style=cs]
namespace ObjectIR.Unity.Runtime
{
    public sealed class UnityIRRuntime { ... }
}
\end{lstlisting}

\code{UnityIRRuntime} is a \kword{sealed} class. It cannot be subclassed.

% ═══════════════════════════════════════════════════════════════════════════════
\section{Inner Types}
% ═══════════════════════════════════════════════════════════════════════════════

\subsection{\texttt{InputFormat} (enum)}

Controls how \code{LoadModule(string, InputFormat)} parses the input string.

\begin{longtable}{lp{9cm}}
  \toprule
  \textbf{Member} & \textbf{Description} \\
  \midrule
  \endhead
  \code{Auto}    & Inspects the first character of the input. A \texttt{\{} or
                   \texttt{[} triggers JSON parsing; everything else is treated as
                   TextIR source. \\[0.3em]
  \code{Json}    & Deserialise as a \typename{ModuleDto} JSON object. \\[0.3em]
  \code{TextIr}  & Parse as human-readable TextIR source via \typename{TextIrParser}
                   and lower to \typename{ModuleDto} through \typename{IrAstLowering}. \\[0.3em]
  \code{Ast}     & Alias for \code{TextIr} --- the string is passed through the
                   same TextIR parser / lowering pipeline. \\[0.3em]
  \code{FobIr}   & The string argument is treated as a \emph{file path} to a FOB/IR
                   v3 binary. The file is read with
                   \code{File.ReadAllBytes} and then decoded by
                   \typename{FobIrReader} + \typename{IrModuleBinaryReader}. \\
  \bottomrule
\end{longtable}

\subsection{\texttt{NativeMethod} (delegate)}

\begin{lstlisting}[style=cs]
public delegate object? NativeMethod(object? instance, object?[] args);
\end{lstlisting}

Signature for all native method callbacks registered in the runtime. \code{instance}
is \code{null} for static calls and a \typename{ManagedObject} wrapping the \code{this}
object for virtual/instance calls. \code{args} contains the positional arguments popped
from the evaluation stack.

\subsection{\texttt{ExceptionSignal} (private)}

An internal \typename{Exception} subclass used as a control-flow sentinel when an IR
\code{throw} instruction fires before a matching \code{try/catch} block can intercept it.
It carries the raw IR exception value in its \code{ExceptionObject} property.

% ═══════════════════════════════════════════════════════════════════════════════
\section{Public Properties}
% ═══════════════════════════════════════════════════════════════════════════════

\begin{longtable}{p{5.5cm}p{2cm}p{6.5cm}}
  \toprule
  \textbf{Property} & \textbf{Type} & \textbf{Description} \\
  \midrule
  \endhead
  \code{OwnerGameObject}               & \typename{GameObject?} &
    The Unity \typename{GameObject} this runtime instance is attached to.
    IR scripts can retrieve a reference to it via
    \code{IRRuntime.GetOwner()} / \code{IRRuntime.GetSelf()}. \\[0.4em]
  \code{EnableReflectionNativeMethods} & \code{bool} &
    When \code{true} (default), any call that does not match a registered
    native binding is retried through CLR reflection before falling back
    to looking up the method in the IR module itself. \\[0.4em]
  \code{OnStep}                        & \code{Action<CallFrame,} \code{InstructionDto>?} &
    Optional step-by-step hook. Invoked \emph{before} each instruction
    is executed, giving host code a way to implement debuggers or
    instruction-level tracing. \\[0.4em]
  \code{OnException}                   & \code{Action<Exception>?} &
    Invoked whenever an exception escapes the IR execution loop, either
    from an unhandled IR \code{throw} or from a CLR exception thrown
    inside a native binding. \\[0.4em]
  \code{CallStack}                     & \code{IReadOnlyList<CallFrame>} &
    Read-only view of the current call stack, ordered from outermost
    (index 0) to innermost frame. \\[0.4em]
  \code{StaticFields}                  & \code{IReadOnlyDictionary} \code{<(string,string),} \code{object?>} &
    All static field values that have been written by executing IR
    \code{stsfld} instructions, keyed by \texttt{(declaringType, fieldName)}. \\
  \bottomrule
\end{longtable}

% ═══════════════════════════════════════════════════════════════════════════════
\section{Public API --- Methods}
% ═══════════════════════════════════════════════════════════════════════════════

\subsection{Constructor}

\begin{lstlisting}[style=cs]
public UnityIRRuntime(GameObject? owner = null)
\end{lstlisting}

Creates a new runtime instance. On first call per process the constructor triggers
\emph{one-time, thread-safe} registration of all exported \code{UnityEngine} types
via \code{EnsureUnityTypesRegistered()} and then calls \code{RegisterManualBindings()}
to set up the hand-written fast-path native methods.

\subsection{Module Loading}

\subsubsection{\texttt{LoadModule(string, InputFormat)}}
\begin{lstlisting}[style=cs]
public void LoadModule(string input, InputFormat format = InputFormat.Auto)
\end{lstlisting}
Loads (or replaces) the active IR module from a string. Format detection is
performed automatically when \code{format} is \code{Auto}.

\subsubsection{\texttt{LoadModule(ModuleNode)}}
\begin{lstlisting}[style=cs]
public void LoadModule(ModuleNode astModule)
\end{lstlisting}
Loads a module from an already-parsed \typename{ModuleNode} AST. The node is lowered
via \typename{IrAstLowering} into a \typename{ModuleDto} before execution.

\subsubsection{\texttt{LoadModule(byte[])}}
\begin{lstlisting}[style=cs]
public void LoadModule(byte[] fobBytes)
\end{lstlisting}
Loads a module from raw FOB/IR v3 binary bytes. Internally calls
\typename{FobIrReader} to strip the FOB envelope and then
\typename{IrModuleBinaryReader} to decode the payload.

\subsubsection{\texttt{LoadFobFile(string)}}
\begin{lstlisting}[style=cs]
public void LoadFobFile(string path)
\end{lstlisting}
Convenience wrapper: reads the file at \code{path} and delegates to
\code{LoadModule(byte[])}.

\subsubsection{\texttt{LoadProgram(string, InputFormat)}}
\begin{lstlisting}[style=cs]
public void LoadProgram(string input, InputFormat format = InputFormat.Auto)
\end{lstlisting}
Like \code{LoadModule} but also clears the call stack and all static field
state, effectively doing a full reset before replacing the module.

\subsection{Method Invocation}

\subsubsection{\texttt{CallMethod(string, object?[])}}
\begin{lstlisting}[style=cs]
public object? CallMethod(string qualifiedName, params object?[] args)
\end{lstlisting}
Locates and executes a method by qualified name (\texttt{"TypeName.MethodName"} or
just \texttt{"MethodName"}). Returns the topmost value on the frame's evaluation stack
after execution, or \code{null} for void methods.

Throws \typename{InvalidOperationException} when no module is loaded or the type
is not found; throws \typename{MissingMethodException} when the method cannot be
located.

\subsubsection{\texttt{CallMethod<TResult>(string, object?[])}}
\begin{lstlisting}[style=cs]
public TResult CallMethod<TResult>(string qualifiedName, params object?[] args)
\end{lstlisting}
Generic overload that casts (or converts via \typename{ValueHelpers}\code{.ConvertTo})
the return value to \code{TResult}. Returns \code{default!} when the result is
\code{null}.

\subsubsection{\texttt{Run()}}
\begin{lstlisting}[style=cs]
public void Run()
\end{lstlisting}
Executes the canonical entry point: looks up \code{Program.Main} (static, no
parameters) in the loaded module and runs it.

\subsection{Module Inspection}

\subsubsection{\texttt{HasMethod(string)}}
\begin{lstlisting}[style=cs]
public bool HasMethod(string qualifiedName)
\end{lstlisting}
Returns \code{true} when the loaded module contains a method matching
\code{qualifiedName} (\texttt{"TypeName.MethodName"} or just \texttt{"MethodName"}).
This performs IR metadata inspection only --- no code is executed.

\subsection{Static Field Access}

\subsubsection{\texttt{GetStaticField(string, string)}}
\begin{lstlisting}[style=cs]
public object? GetStaticField(string declaringType, string fieldName)
\end{lstlisting}
Reads a static field value from the runtime's internal static field store.

\subsubsection{\texttt{GetStaticField(string)}}
\begin{lstlisting}[style=cs]
public object? GetStaticField(string qualifiedName)
\end{lstlisting}
Overload that accepts \texttt{"TypeName.FieldName"} as a single string.

\subsection{CLR Type \& Method Registration}

\subsubsection{\texttt{RegisterClrType(string, Type)}}
\begin{lstlisting}[style=cs]
public static void RegisterClrType(string irTypeName, Type clrType)
\end{lstlisting}
Registers a CLR type under an IR type name so that the runtime's reflection
resolution layer can find it. The registration is \emph{process-global} (static).

\subsubsection{\texttt{RegisterClrType<T>(string?)}}
\begin{lstlisting}[style=cs]
public static void RegisterClrType<T>(string? irTypeName = null)
\end{lstlisting}
Generic convenience wrapper. If \code{irTypeName} is \code{null} the CLR full name
of \code{T} is used as the IR type name.

\subsubsection{\texttt{RegisterNativeMethod(string, NativeMethod)}}
\begin{lstlisting}[style=cs]
public void RegisterNativeMethod(string signature, NativeMethod method)
\end{lstlisting}
Adds or replaces a native method binding in the \emph{instance-level}
\code{\_nativeMethods} dictionary. The \code{signature} format is:

\begin{center}
  \code{DeclaringType.MethodName(param1Type,param2Type,...)}
\end{center}

Type names in signatures are \emph{case-sensitive} but the dispatcher also tries a
lower-cased, short-name fallback so that \code{(float32)} and \code{(single)} both
resolve.

% ═══════════════════════════════════════════════════════════════════════════════
\section{VM Instruction Set}
% ═══════════════════════════════════════════════════════════════════════════════

The execution engine is a simple stack machine. The call stack is a
\code{Stack<CallFrame>} and each frame has its own evaluation stack
(\code{Stack<object?>}). All values are untyped (\code{object?}) and conversions are
performed lazily by \typename{ValueHelpers}.

\subsection{Stack / Flow Control Instructions}

\begin{longtable}{lp{10cm}}
  \toprule
  \textbf{Opcode} & \textbf{Semantics} \\
  \midrule
  \endhead
  \code{nop}      & No operation. \\
  \code{dup}      & Duplicates the top-of-stack value. \\
  \code{pop}      & Discards the top-of-stack value. \\
  \code{ldnull}   & Pushes \code{null} onto the stack. \\
  \code{ret}      & Returns from the current call frame; pops a return value from the
                    evaluation stack (if present) and pushes it onto the caller's stack. \\
  \code{throw}    & Pops the top-of-stack as an exception object and stores it in the
                    current frame's \code{PendingException}. \\
  \code{break}    & Raises \typename{BreakSignal} to exit the nearest \code{while} body. \\
  \code{continue} & Raises \typename{ContinueSignal} to jump to the next iteration of
                    the nearest \code{while} body. \\
  \bottomrule
\end{longtable}

\subsection{Load / Store Instructions}

\begin{longtable}{lp{10cm}}
  \toprule
  \textbf{Opcode} & \textbf{Semantics} \\
  \midrule
  \endhead
  \code{ldc}    & Load constant. Operand: \code{\{value, type\}}. Value is converted
                  to the target type via \typename{ValueHelpers}\code{.ConvertTo}. \\
  \code{ldstr}  & Load string constant. Operand: \code{\{value\}}. \\
  \code{ldarg}  & Load argument by index (\code{operand.index}) or name
                  (\code{operand.argumentName}). The special name \texttt{"this"}
                  loads the \code{CallFrame.This} reference. \\
  \code{starg}  & Store top-of-stack into a named argument slot. \\
  \code{ldloc}  & Load local variable by name (\code{operand.localName}). \\
  \code{stloc}  & Store top-of-stack into a named local variable. \\
  \code{ldfld}  & Load instance field. Pops the \typename{ManagedObject} instance
                  (or uses \code{frame.This} if none is on the stack). Tries the
                  embedded CLR instance first via reflection; falls back to the
                  \typename{ManagedObject} field dictionary. \\
  \code{stfld}  & Store top-of-stack into an instance field. Mirror of \code{ldfld}. \\
  \code{ldsfld} & Load a static field. Operand: \code{\{field: \{declaringType, name\}\}}.
                  Attempts CLR property/field access first; falls back to the
                  runtime's \code{\_staticFields} dictionary. \\
  \code{stsfld} & Store top-of-stack into a static field. Mirror of \code{ldsfld}. \\
  \code{ldelem} & Pop index then array (\code{List<object?>}) and push the element. \\
  \code{stelem} & Pop value, index, then array and write the element (grows the list
                  with \code{null} fill if needed). \\
  \bottomrule
\end{longtable}

\subsection{Arithmetic \& Comparison Instructions}

\begin{longtable}{lp{10cm}}
  \toprule
  \textbf{Opcode} & \textbf{Semantics} \\
  \midrule
  \endhead
  \code{add}  & Pop \(b\), pop \(a\), push \(a + b\). \\
  \code{sub}  & Pop \(b\), pop \(a\), push \(a - b\). \\
  \code{mul}  & Pop \(b\), pop \(a\), push \(a \times b\). \\
  \code{div}  & Pop \(b\), pop \(a\), push \(a / b\). \\
  \code{rem}  & Pop \(b\), pop \(a\), push \(a \bmod b\). \\
  \code{neg}  & Pop \(v\), push \(-v\). \\
  \code{not}  & Pop \(v\), push \(\neg v\) (boolean negation). \\
  \code{ceq}  & Push \code{true} iff \(a = b\). \\
  \code{cne}  & Push \code{true} iff \(a \neq b\). \\
  \code{cgt}  & Push \code{true} iff \(a > b\). \\
  \code{cge}  & Push \code{true} iff \(a \geq b\). \\
  \code{clt}  & Push \code{true} iff \(a < b\). \\
  \code{cle}  & Push \code{true} iff \(a \leq b\). \\
  \bottomrule
\end{longtable}

Arithmetic selects integer (\code{int64}) or floating-point (\code{double}) mode
based on the runtime types of \(a\) and \(b\). If either operand is a \code{float},
\code{double}, or \code{string}, double mode is used; otherwise integer mode.

String comparisons use \code{StringComparison.Ordinal}. Boolean comparisons only
support \code{ceq} / \code{cne}.

\subsection{Type Conversion Instructions}

\begin{longtable}{lp{10cm}}
  \toprule
  \textbf{Opcode} & \textbf{Semantics} \\
  \midrule
  \endhead
  \code{conv}       & Pop \(v\) and push \code{ValueHelpers.ConvertTo(targetType, v)}.
                      Operand: \code{\{targetType\}}. \\
  \code{castclass}  & Asserts that the top-of-stack \typename{ManagedObject} has a
                      matching \code{TypeName}. Throws \typename{InvalidCastException}
                      on failure. Passes \code{null} through unchanged. \\
  \code{isinst}     & Pushes \code{true}/\code{false} depending on whether the
                      top-of-stack is a \typename{ManagedObject} with the expected
                      \code{TypeName}. \\
  \bottomrule
\end{longtable}

\subsection{Object Creation Instructions}

\begin{longtable}{lp{10cm}}
  \toprule
  \textbf{Opcode} & \textbf{Semantics} \\
  \midrule
  \endhead
  \code{newobj} & Creates a new \typename{ManagedObject} with the given type name.
                  If the type is registered in \code{s\_registeredClrTypes} a CLR
                  instance is created via \code{Activator.CreateInstance} and
                  embedded as \code{"clrInstance"} metadata. IR-defined attributes
                  are also attached. Operand: \code{\{type\}}. \\
  \code{newarr} & Pushes an empty \code{List<object?>} as a dynamically-typed array.
                  Operand: \code{\{elementType\}} (ignored at runtime; stored for
                  source-level tooling). \\
  \bottomrule
\end{longtable}

\subsection{Call Instructions}

\begin{longtable}{lp{10cm}}
  \toprule
  \textbf{Opcode} & \textbf{Semantics} \\
  \midrule
  \endhead
  \code{call}     & Static / non-virtual call. Pops arguments right-to-left, looks up
                    the native binding, reflection, or IR method, pushes the result
                    (if non-void). Operand: \code{\{method: CallTargetDto\}}. \\
  \code{callvirt} & Virtual / instance call. Same as \code{call} but also pops the
                    \typename{ManagedObject} instance from the stack and passes it as
                    the \code{instance} parameter to the native binding. \\
  \bottomrule
\end{longtable}

Method dispatch order for both \code{call} and \code{callvirt}:

\begin{enumerate}[leftmargin=2em]
  \item Exact signature lookup in \code{\_nativeMethods}.
  \item Simplified-parameter (lowercased, short-name) fallback in
        \code{\_nativeMethods}.
  \item CLR reflection via \code{TryInvokeReflectedNative} (when
        \code{EnableReflectionNativeMethods} is \code{true}).
  \item IR method lookup in the loaded module.
\end{enumerate}

\subsection{Control-Flow Instructions}

\begin{longtable}{lp{10cm}}
  \toprule
  \textbf{Opcode} & \textbf{Semantics} \\
  \midrule
  \endhead
  \code{if}    & Evaluates the embedded \typename{ConditionDto} and executes either
                 \code{thenBlock} or \code{elseBlock}. Operand:
                 \code{\{condition, thenBlock, elseBlock?\}}. \\
  \code{while} & Evaluates the condition before each iteration; executes \code{body}.
                 Honours \typename{BreakSignal} / \typename{ContinueSignal} for
                 \code{break}/\code{continue} support. \\
  \code{try}   & Structured exception handling. Contains \code{tryBlock},
                 \code{catchBlocks[]}, and an optional \code{finallyBlock}. Each
                 catch block may specify an \code{exceptionType} (IR type name or CLR
                 type name) or catch all exceptions. The \code{finally} block runs
                 unconditionally after catch processing. \\
  \bottomrule
\end{longtable}

\subsubsection{Condition System}

Conditions (\typename{ConditionDto}) are evaluated by \code{EvaluateCondition} and
support four strategies:

\begin{longtable}{lp{10cm}}
  \toprule
  \textbf{Kind} & \textbf{Semantics} \\
  \midrule
  \endhead
  \code{stack}      & Pops the top-of-stack and converts to bool via
                      \typename{ValueHelpers}\code{.ToBool}. \\
  \code{binary}     & Pops right then left, applies \code{operation} (one of the
                      comparison opcodes: \code{ceq}, \code{cne}, \ldots). \\
  \code{expression} & Executes a single embedded \typename{InstructionDto} then pops
                      and converts to bool. \\
  \code{block}      & Executes an embedded \typename{InstructionDto}\code{[]} block
                      then pops and converts to bool. \\
  \bottomrule
\end{longtable}

% ═══════════════════════════════════════════════════════════════════════════════
\section{Data-Transfer Object Types}
% ═══════════════════════════════════════════════════════════════════════════════

These types live in \texttt{UnityIRRuntime.Types.cs} and are the in-memory
representation of a compiled IR module.

\subsection{\texttt{ModuleDto}}

\begin{lstlisting}[style=cs]
public sealed class ModuleDto
{
    public string      name    { get; set; }
    public string      version { get; set; }
    public JsonElement metadata  { get; set; }
    public TypeDto[]   types     { get; set; }
    public JsonElement functions { get; set; }
}
\end{lstlisting}

Top-level container for an IR module. \code{types} holds all class and interface
definitions. \code{functions} is reserved for future module-level free functions
(currently unused by the VM).

\subsection{\texttt{TypeDto}}

\begin{lstlisting}[style=cs]
public sealed class TypeDto
{
    public string  kind; public string name; public string? _namespace;
    public string? access; public bool isAbstract; public bool isSealed;
    public string? baseType;
    public AttributeDto[]  attributes;
    public FieldDto[]      fields;
    public MethodDto[]     methods;
    public string[]        interfaces;
    public JsonElement baseInterfaces, genericParameters, properties;
}
\end{lstlisting}

Represents one class or interface. \code{kind} is \texttt{"class"} or
\texttt{"interface"}.

\subsection{\texttt{MethodDto}}

\begin{lstlisting}[style=cs]
public sealed class MethodDto
{
    public string  name, returnType;
    public string? access;
    public bool isStatic, isVirtual, isOverride, isAbstract, isConstructor;
    public AttributeDto[]      attributes;
    public ParameterDto[]      parameters;
    public LocalVariableDto[]  localVariables;
    public int                 instructionCount;
    public InstructionDto[]    instructions;
}
\end{lstlisting}

\subsection{\texttt{InstructionDto}}

\begin{lstlisting}[style=cs]
public sealed class InstructionDto
{
    public string      opCode  { get; set; }
    public JsonElement operand { get; set; }
}
\end{lstlisting}

The operand is stored as a raw \typename{JsonElement} to allow the VM to lazily
deserialise only the fields it needs for each opcode.

\subsection{\texttt{CallTargetDto}}

\begin{lstlisting}[style=cs]
public sealed class CallTargetDto
{
    public string   declaringType;
    public string   name;
    public string   returnType;
    public string[] parameterTypes;
}
\end{lstlisting}

Encodes the target of a \code{call} or \code{callvirt} instruction. The VM builds
the lookup key as:

\begin{center}
  \code{declaringType.name(parameterTypes[0],parameterTypes[1],...)}
\end{center}

\subsection{\texttt{ConditionDto}}

\begin{lstlisting}[style=cs]
public sealed class ConditionDto
{
    public string           kind;
    public string?          operation;
    public InstructionDto?  expression;
    public InstructionDto[]? block;
    // [JsonIgnore]
    public ConditionKind Kind { get; }   // computed from 'kind'
}
\end{lstlisting}

\subsection{\texttt{FieldDto}, \texttt{ParameterDto}, \texttt{LocalVariableDto}, \texttt{AttributeDto}, \texttt{FieldRefDto}}

Supporting DTOs that describe field declarations, method parameters, local
variables, custom attributes, and field references respectively.

% ═══════════════════════════════════════════════════════════════════════════════
\section{Runtime Support Types}
% ═══════════════════════════════════════════════════════════════════════════════

\subsection{\texttt{CallFrame}}

Represents a single activation record on the call stack.

\begin{longtable}{lp{9cm}}
  \toprule
  \textbf{Member} & \textbf{Description} \\
  \midrule
  \endhead
  \code{Method}          & The \typename{MethodDto} being executed. \\
  \code{IP}              & Instruction pointer (index into \code{Method.instructions}). \\
  \code{EvalStack}       & Per-frame evaluation stack (\code{Stack<object?>}). \\
  \code{Args}            & Array of argument values for this call. \\
  \code{ArgNameToIndex}  & Maps parameter names to their index in \code{Args}. \\
  \code{Locals}          & Dictionary mapping local variable names to their values.
                           Pre-populated with \code{null} for each declared local. \\
  \code{This}            & The object instance for instance/virtual calls. \\
  \code{PendingException}& Set by \code{throw}; propagates outward until a matching
                           \code{catch} block consumes it. \\
  \bottomrule
\end{longtable}

\subsection{\texttt{ManagedObject}}

Represents any heap-allocated value in the IR runtime (classes, Unity value types,
exception objects, \ldots).

\begin{longtable}{lp{9cm}}
  \toprule
  \textbf{Member} & \textbf{Description} \\
  \midrule
  \endhead
  \code{TypeName}                     & The IR type name this object was instantiated as. \\
  \code{GetField / SetField}          & Access named field slots (string-keyed dictionary). \\
  \code{GetMetadata / SetMetadata}    & Side-channel key/value store used to attach CLR
                                        objects (\code{"clrInstance"}) or other host data
                                        without polluting the IR field namespace. \\
  \code{AddAttribute / GetAttribute<T>} & Attach and retrieve attribute instances created
                                          from IR \typename{AttributeDto} declarations. \\
  \code{AttachAttributesFromIR}       & Instantiates all attributes declared on a
                                        \typename{TypeDto} and attaches them. \\
  \bottomrule
\end{longtable}

\textbf{Unity struct boxing.} Unity value types (\typename{Vector3},
\typename{Vector2}, \typename{Quaternion}, \typename{Color}) are represented as
\typename{ManagedObject} instances with:

\begin{itemize}[leftmargin=2em]
  \item \code{TypeName} set to \texttt{"UnityEngine.Vector3"} (etc.)
  \item \code{"clrInstance"} metadata holding the actual CLR struct.
  \item Individual component fields (\code{x}, \code{y}, \code{z}, etc.) mirrored
        as IR-readable field slots for scripts that access them directly.
\end{itemize}

\subsection{\texttt{ValueHelpers} (internal)}

Stateless helper class providing universal type coercion.

\begin{longtable}{lp{9cm}}
  \toprule
  \textbf{Method} & \textbf{Description} \\
  \midrule
  \endhead
  \code{ToBool(object?)}             & Converts any value to \code{bool}.
                                       \code{null}/\code{false}/0/empty string → false;
                                       everything else → true. \\
  \code{ToDouble(object?)}           & Converts any numeric / string value to
                                       \code{double}. \\
  \code{ToInt64(object?)}            & Converts any numeric / string value to
                                       \code{long}. \\
  \code{ConvertTo(string, object?)}  & Converts to the type identified by an IR
                                       type name string (e.g.\ \texttt{"int32"},
                                       \texttt{"float32"}, \texttt{"bool"}). \\
  \code{ConvertTo(Type, object?)}    & CLR-type-aware overload. Handles all primitive
                                       types, enums, and falls back to
                                       \code{Convert.ChangeType}. \\
  \bottomrule
\end{longtable}

\subsection{\texttt{BreakSignal} / \texttt{ContinueSignal} (internal)}

Lightweight \typename{Exception} subclasses used as structured control-flow tokens
for \code{break} and \code{continue} within \code{while} loops. They are caught
exclusively inside the \code{while} instruction handler and never escape to user code.

% ═══════════════════════════════════════════════════════════════════════════════
\section{Type Resolution \& CLR Integration}
% ═══════════════════════════════════════════════════════════════════════════════

\subsection{Type Name Resolution Pipeline}

\code{ResolveClrType(string typeName)} applies the following steps in order:

\begin{enumerate}[leftmargin=2em]
  \item Exact lookup in \code{s\_registeredClrTypes} (process-global registry).
  \item \code{Type.GetType(typeName)} against mscorlib.
  \item \code{MapTypeNameToClrType} — maps IR primitive names to CLR types
        (see table below).
  \item Short-name lookup in \code{s\_registeredClrTypes} (strips leading
        namespace).
  \item Full scan of all loaded assemblies via
        \code{AppDomain.CurrentDomain.GetAssemblies()}.
\end{enumerate}

\subsection{IR Primitive Type Name Mapping}

\begin{longtable}{ll}
  \toprule
  \textbf{IR type name(s)} & \textbf{CLR type} \\
  \midrule
  \endhead
  \code{void}, \code{system.void}           & \code{void} \\
  \code{bool}, \code{system.boolean}        & \code{bool} \\
  \code{int8}, \code{system.sbyte}          & \code{sbyte} \\
  \code{uint8}, \code{system.byte}          & \code{byte} \\
  \code{int16}, \code{system.int16}         & \code{short} \\
  \code{uint16}, \code{system.uint16}       & \code{ushort} \\
  \code{int32}, \code{system.int32}         & \code{int} \\
  \code{uint32}, \code{system.uint32}       & \code{uint} \\
  \code{int64}, \code{system.int64}         & \code{long} \\
  \code{uint64}, \code{system.uint64}       & \code{ulong} \\
  \code{float32}, \code{single}, \code{system.single} & \code{float} \\
  \code{float64}, \code{double}, \code{system.double} & \code{double} \\
  \code{char}, \code{system.char}           & \code{char} \\
  \code{string}, \code{system.string}       & \code{string} \\
  \code{object}, \code{system.object}       & \code{object} \\
  \code{decimal}, \code{system.decimal}     & \code{decimal} \\
  \code{datetime}, \code{system.datetime}   & \code{DateTime} \\
  \code{timespan}, \code{system.timespan}   & \code{TimeSpan} \\
  \code{guid}, \code{system.guid}           & \code{Guid} \\
  \code{system.console}                     & \code{Console} \\
  \bottomrule
\end{longtable}

\subsection{Unity Type Auto-Registration}

At first runtime instantiation \code{RegisterUnityClrTypes()} runs once (under a
static lock) and:

\begin{enumerate}[leftmargin=2em]
  \item Explicitly registers the \(\approx 70\) most common Unity types under
        both their fully-qualified names (e.g.\ \texttt{"UnityEngine.Debug"}) and
        their short names (e.g.\ \texttt{"Debug"}).
  \item Calls \code{AutoRegisterAssemblyTypes} on
        \texttt{UnityEngine.CoreModule} and
        \texttt{UnityEngine.SceneManagement} to bulk-register every exported public
        type from those assemblies.
\end{enumerate}

Explicit registrations listed below by category:

\begin{center}
\begin{tabular}{ll}
  \toprule
  \textbf{Category} & \textbf{Types} \\
  \midrule
  Core objects      & \code{Debug}, \code{GameObject}, \code{Transform},
                      \code{Component}, \code{Object},
                      \code{MonoBehaviour}, \code{Behaviour} \\
  Math / geometry   & \code{Vector2/3/4}, \code{Vector2Int/3Int},
                      \code{Quaternion}, \code{Color}, \code{Color32},
                      \code{Mathf}, \code{Bounds}, \code{Rect},
                      \code{Ray}, \code{Ray2D},
                      \code{RaycastHit}, \code{RaycastHit2D},
                      \code{LayerMask}, \code{Matrix4x4} \\
  Engine services   & \code{Time}, \code{Input}, \code{Screen},
                      \code{Application}, \code{Cursor}, \code{Random} \\
  Physics           & \code{Physics}, \code{Physics2D}, \code{Collider},
                      \code{Collider2D}, \code{Rigidbody}, \code{Rigidbody2D},
                      \code{CharacterController} \\
  Rendering         & \code{Camera}, \code{Light}, \code{Renderer},
                      \code{MeshRenderer}, \code{SpriteRenderer},
                      \code{Material}, \code{Texture2D}, \code{Sprite} \\
  Audio             & \code{AudioSource}, \code{AudioClip},
                      \code{AudioListener} \\
  Animation         & \code{Animator}, \code{Animation} \\
  Scenes            & \code{SceneManager} \\
  \bottomrule
\end{tabular}
\end{center}

% ═══════════════════════════════════════════════════════════════════════════════
\section{Native Method Bindings}
% ═══════════════════════════════════════════════════════════════════════════════

Native bindings are stored in the \code{\_nativeMethods} dictionary and are
looked up by \emph{call signature}. The following subsections list every
hand-written binding registered by \code{RegisterManualBindings()}.
All bindings are registered for both the fully-qualified form (e.g.\
\code{UnityEngine.Debug}) and the short form (\code{Debug}).

\subsection{System.Console (redirected to Debug.Log)}

All \code{System.Console.WriteLine} / \code{Write} variants redirect to
\code{Debug.Log} so that scripts ported from a console environment produce
visible output in the Unity Editor. \code{ReadLine}, \code{Clear}, and
\code{Beep} are silently no-ops.

\subsection{UnityEngine.Debug}

Full coverage of \code{Log}, \code{LogWarning}, \code{LogError},
\code{LogException}, \code{Assert} (with and without message), \code{DrawLine},
\code{DrawRay}, \code{Break}, \code{ClearDeveloperConsole}, and the
\code{isDebugBuild} property.

\subsection{UnityEngine.Vector3}

\begin{itemize}[leftmargin=2em, noitemsep]
  \item Constructor: \code{(float32, float32, float32)} and \code{(float32, float32)}.
  \item Operators: \code{+}, \code{-}, \code{*} (scalar), \code{/} (scalar),
        unary negation, \code{==}, \code{!=}.
  \item Static methods: \code{Distance}, \code{Dot}, \code{Cross},
        \code{Lerp}, \code{LerpUnclamped}, \code{MoveTowards}, \code{Normalize},
        \code{Scale}, \code{Reflect}, \code{Angle}, \code{SignedAngle},
        \code{ClampMagnitude}, \code{Project}, \code{ProjectOnPlane},
        \code{RotateTowards}.
  \item Instance properties: \code{normalized}, \code{magnitude}, \code{sqrMagnitude}.
  \item Static constants: \code{zero}, \code{one}, \code{up}, \code{down},
        \code{left}, \code{right}, \code{forward}, \code{back},
        \code{positiveInfinity}, \code{negativeInfinity}.
\end{itemize}

\subsection{UnityEngine.Vector2}

Constructor, arithmetic operators (\code{+}, \code{-}, \code{*}, \code{/}),
\code{Distance}, \code{Dot}, \code{Lerp}, \code{MoveTowards}, \code{Reflect},
\code{Angle}, instance properties \code{normalized} / \code{magnitude}, and static
constants \code{zero}, \code{one}, \code{up}, \code{down}, \code{left},
\code{right}.

\subsection{UnityEngine.Quaternion}

Constructor, \code{Euler}, \code{AngleAxis}, \code{LookRotation},
\code{Slerp}, \code{Lerp}, \code{RotateTowards}, \code{Angle},
\code{Inverse}, \code{FromToRotation}, multiplication operators
(\code{Quaternion}$\times$\code{Quaternion} and \code{Quaternion}$\times$\code{Vector3}),
property \code{eulerAngles}, and constant \code{identity}.

\subsection{UnityEngine.Color}

Constructors with 3 and 4 float components; \code{Lerp}, \code{LerpUnclamped};
static colour presets: \code{red}, \code{green}, \code{blue}, \code{white},
\code{black}, \code{yellow}, \code{cyan}, \code{magenta}, \code{clear},
\code{gray}, \code{grey}.

\subsection{UnityEngine.Time}

\code{deltaTime}, \code{fixedDeltaTime} (get + set), \code{time},
\code{unscaledTime}, \code{unscaledDeltaTime}, \code{timeScale} (get + set),
\code{frameCount}, \code{realtimeSinceStartup}, \code{smoothDeltaTime},
\code{timeSinceLevelLoad}, \code{renderedFrameCount}.

\subsection{UnityEngine.Input}

\code{GetKey}, \code{GetKeyDown}, \code{GetKeyUp},
\code{GetAxis}, \code{GetAxisRaw},
\code{GetButton}, \code{GetButtonDown}, \code{GetButtonUp},
\code{GetMouseButton}, \code{GetMouseButtonDown}, \code{GetMouseButtonUp},
\code{GetTouch}, properties \code{mousePosition}, \code{mouseScrollDelta},
\code{anyKey}, \code{anyKeyDown}, \code{inputString}, \code{touchCount},
\code{deviceOrientation}, \code{acceleration}.

\subsection{UnityEngine.Mathf}

Complete coverage of mathematical functions:

\begin{itemize}[leftmargin=2em, noitemsep]
  \item Trigonometry: \code{Sin}, \code{Cos}, \code{Tan}, \code{Asin},
        \code{Acos}, \code{Atan}, \code{Atan2}.
  \item Exponential: \code{Sqrt}, \code{Pow}, \code{Log}, \code{Log10},
        \code{Exp}.
  \item Rounding: \code{Floor}, \code{FloorToInt}, \code{Ceil},
        \code{CeilToInt}, \code{Round}, \code{RoundToInt}.
  \item Clamping / interpolation: \code{Clamp}, \code{Clamp01},
        \code{Lerp}, \code{LerpUnclamped}, \code{LerpAngle},
        \code{InverseLerp}, \code{SmoothStep}, \code{SmoothDamp}.
  \item Utility: \code{Abs}, \code{Min}, \code{Max}, \code{Sign},
        \code{PingPong}, \code{Repeat}, \code{DeltaAngle},
        \code{MoveTowards}, \code{MoveTowardsAngle}, \code{Approximately},
        \code{IsPowerOfTwo}, \code{NextPowerOfTwo}.
  \item Constants: \code{PI}, \code{Infinity}, \code{NegativeInfinity},
        \code{Epsilon}, \code{Deg2Rad}, \code{Rad2Deg}.
\end{itemize}

\subsection{UnityEngine.Random}

\code{Range(float,float)}, \code{Range(int,int)}, \code{value},
\code{insideUnitCircle}, \code{insideUnitSphere}, \code{onUnitSphere},
\code{rotation}, \code{rotationUniform}.

\subsection{UnityEngine.Application}

\code{Quit()}, \code{Quit(int32)}, \code{OpenURL}; read-only properties
\code{isPlaying}, \code{isEditor}, \code{isFocused}, \code{platform},
\code{dataPath}, \code{persistentDataPath}, \code{temporaryCachePath},
\code{streamingAssetsPath}, \code{version}, \code{productName},
\code{companyName}, \code{unityVersion}; read-write
\code{targetFrameRate}, \code{runInBackground}.

\subsection{UnityEngine.SceneManagement.SceneManager}

\code{LoadScene(string)}, \code{LoadScene(int32)},
\code{LoadScene(string, LoadSceneMode)},
\code{UnloadSceneAsync(string)},
\code{GetActiveScene()} (returns active scene name as string),
\code{GetActiveSceneName()}, \code{GetSceneByName(string)},
\code{sceneCount}.

\subsection{UnityEngine.Screen}

\code{width}, \code{height}, \code{dpi}, \code{fullScreen} (get + set),
\code{currentResolution}, \code{SetResolution}.

\subsection{UnityEngine.Cursor}

\code{visible} (get + set), \code{lockState} (get + set as \code{int32}).

\subsection{UnityEngine.Camera}

\code{main}, \code{current}; instance methods \code{ScreenToWorldPoint},
\code{WorldToScreenPoint}, \code{ScreenPointToRay}; properties
\code{fieldOfView} (get + set), \code{nearClipPlane}, \code{farClipPlane},
\code{orthographic} (get + set), \code{backgroundColor} (get + set).

\subsection{UnityEngine.Physics}

\code{Raycast(vector3, vector3)},
\code{Raycast(vector3, vector3, float32)},
\code{OverlapSphere}, \code{gravity} (get + set),
\code{bounceThreshold}.

\subsection{UnityEngine.Physics2D}

\code{Raycast(vector2, vector2)},
\code{Raycast(vector2, vector2, float32)},
\code{OverlapCircle}, \code{gravity} (get + set).

\subsection{UnityEngine.GameObject}

Static: \code{Find}, \code{FindWithTag}, \code{FindGameObjectWithTag},
\code{CreatePrimitive};
constructors: \code{GameObject(string)} and \code{GameObject()}.

Instance: \code{name} (get + set), \code{tag} (get + set),
\code{activeSelf}, \code{activeInHierarchy}, \code{layer} (get + set),
\code{transform}, \code{SetActive}, \code{CompareTag}, \code{AddComponent},
\code{GetComponent}, \code{GetComponentInChildren}, \code{GetComponentInParent},
\code{SendMessage} (with and without argument).

\subsection{UnityEngine.Object}

\code{Instantiate} (3 overloads: plain, with position and rotation, with parent),
\code{Destroy} (with and without delay), \code{DestroyImmediate},
\code{DontDestroyOnLoad}, \code{FindObjectOfType},
\code{name} (get + set).

\subsection{UnityEngine.Transform}

\begin{itemize}[leftmargin=2em, noitemsep]
  \item Position / rotation / scale properties:
        \code{position}, \code{localPosition}, \code{rotation},
        \code{localRotation}, \code{localScale}, \code{lossyScale},
        \code{eulerAngles}, \code{localEulerAngles} (all get + set).
  \item Direction helpers (read-only): \code{forward}, \code{right}, \code{up}.
  \item Hierarchy: \code{parent} (get + set), \code{childCount},
        \code{gameObject}, \code{GetChild}, \code{GetChildCount},
        \code{Find}, \code{IsChildOf},
        \code{SetParent} (with and without \code{worldPositionStays}),
        \code{SetAsFirstSibling}, \code{SetAsLastSibling},
        \code{DetachChildren}.
  \item Movement / rotation methods:
        \code{Translate}, \code{Rotate}, \code{RotateAround},
        \code{LookAt} (by direction vector or target Transform),
        \code{SetPositionAndRotation}.
  \item Coordinate space conversions:
        \code{TransformPoint}, \code{InverseTransformPoint},
        \code{TransformDirection}, \code{InverseTransformDirection}.
\end{itemize}

\subsection{Owner / Self Access Bindings}

The following four signatures allow IR scripts to obtain a reference to the
runtime's \code{OwnerGameObject} without requiring a method argument:

\begin{lstlisting}[style=cs]
IRRuntime.GetOwner()     -> ManagedObject (GameObject)
IRRuntime.GetSelf()      -> ManagedObject (GameObject)
ObjectIR.GetOwner()      -> ManagedObject (GameObject)
ObjectIR.GetSelf()       -> ManagedObject (GameObject)
ObjectIR.GetGameObject() -> ManagedObject (GameObject)
\end{lstlisting}

% ═══════════════════════════════════════════════════════════════════════════════
\section{Module Format Support}
% ═══════════════════════════════════════════════════════════════════════════════

\subsection{JSON (\texttt{ModuleDto})}

Standard JSON serialisation matching the \typename{ModuleDto} schema (camelCase
property names). Deserialised with \code{System.Text.Json.JsonSerializer}.

\subsection{TextIR}

Human-readable ObjectIR source text. Parsed by \typename{TextIrParser} from the
\code{ObjectIR.AST} assembly, then lowered to \typename{ModuleDto} by the embedded
\typename{IrAstLowering} class. \typename{IrAstLowering} handles:

\begin{itemize}[leftmargin=2em, noitemsep]
  \item \code{ModuleNode} → \typename{ModuleDto}
  \item \code{ClassNode} / \code{InterfaceNode} → \typename{TypeDto}
  \item \code{MethodNode} / \code{ConstructorNode} → \typename{MethodDto}
  \item \code{BlockStatement} → \typename{InstructionDto}\code{[]}
        (with \code{if}/\code{while} lowering and condition expression parsing)
  \item Binary condition string parsing (e.g.\ \texttt{"x == 0"} → \code{block}
        condition with \code{ceq} opcode)
\end{itemize}

\subsection{FOB/IR v3 Binary}

A two-stage decode:

\begin{enumerate}[leftmargin=2em]
  \item \typename{FobIrReader} strips the FOB envelope (from the \code{ObjectIR.FobCompiler}
        assembly) and yields the raw payload bytes.
  \item \typename{IrModuleBinaryReader} (embedded in \texttt{UnityIRRuntime.Types.cs})
        reads the payload.
\end{enumerate}

\subsubsection{Binary Format Summary}

\begin{longtable}{lp{9cm}}
  \toprule
  \textbf{Structure} & \textbf{Layout} \\
  \midrule
  \endhead
  String table       & \code{uint32} count; then for each: \code{uint16} byte-length,
                       UTF-8 bytes. All subsequent strings are 2-byte indices into
                       this table. \\
  Module             & name index, version index, \code{uint16} type count, then types. \\
  Type               & kind, name, namespace, access, baseType indices; 1-byte flags
                       (\code{isAbstract}, \code{isSealed}); interface list;
                       field list; method list. \\
  Method             & name, returnType, access indices; 1-byte flags; param list;
                       local list; \code{uint32} instruction count; instructions. \\
  Instruction        & 1-byte opcode; opcode-dependent operands (see below). \\
  \bottomrule
\end{longtable}

\subsubsection{Binary Opcode Map}

\begin{longtable}{llll}
  \toprule
  \textbf{Byte} & \textbf{Opcode} & \textbf{Byte} & \textbf{Opcode} \\
  \midrule
  \endhead
  \code{0x00} & \code{nop}       & \code{0x11} & \code{ceq} \\
  \code{0x01} & \code{dup}       & \code{0x12} & \code{cne} \\
  \code{0x02} & \code{pop}       & \code{0x13} & \code{cgt} \\
  \code{0x03} & \code{ldnull}    & \code{0x14} & \code{cge} \\
  \code{0x0A} & \code{add}       & \code{0x15} & \code{clt} \\
  \code{0x0B} & \code{sub}       & \code{0x16} & \code{cle} \\
  \code{0x0C} & \code{mul}       & \code{0x17} & \code{ldelem} \\
  \code{0x0D} & \code{div}       & \code{0x18} & \code{stelem} \\
  \code{0x0E} & \code{rem}       & \code{0x19} & \code{ret} \\
  \code{0x0F} & \code{neg}       & \code{0x1A} & \code{throw} \\
  \code{0x10} & \code{not}       & \code{0x1B} & \code{break} \\
              &                  & \code{0x1C} & \code{continue} \\
  \midrule
  \code{0x20} & \code{ldstr}     & \code{0x29} & \code{conv} \\
  \code{0x21} & \code{ldarg}     & \code{0x30} & \code{ldc} \\
  \code{0x22} & \code{starg}     & \code{0x31} & \code{ldsfld} \\
  \code{0x23} & \code{ldloc}     & \code{0x32} & \code{stsfld} \\
  \code{0x24} & \code{stloc}     & \code{0x33} & \code{castclass} \\
  \code{0x25} & \code{newobj}    & \code{0x34} & \code{isinst} \\
  \code{0x26} & \code{newarr}    & \code{0x40} & \code{call} \\
  \code{0x27} & \code{ldfld}     & \code{0x41} & \code{callvirt} \\
  \code{0x28} & \code{stfld}     & \code{0x50} & \code{if} \\
  \midrule
  \code{0x51} & \code{while}     & \code{0x52} & \code{try} \\
  \code{0xFF} & \multicolumn{3}{l}{Raw JSON escape (uint32 length prefix + UTF-8 JSON)} \\
  \bottomrule
\end{longtable}

% ═══════════════════════════════════════════════════════════════════════════════
\section{Error Handling}
% ═══════════════════════════════════════════════════════════════════════════════

\subsection{IR Exception Propagation}

\begin{enumerate}[leftmargin=2em]
  \item IR \code{throw} stores the exception object in \code{frame.PendingException}.
  \item The \code{ExecuteBlock} loop stops execution for the current block when
        \code{PendingException} is non-null.
  \item The \code{try} instruction intercepts exceptions; \code{catch} blocks are
        tried in order matching by IR type name or CLR type assignability.
  \item Exceptions that escape all frames are passed to \code{OnException} (if
        set) and re-thrown as a CLR \typename{Exception}.
\end{enumerate}

\subsection{IR Exception to CLR Exception Conversion}

When an IR \typename{ManagedObject} exception escapes unhandled,
\code{BuildExceptionFromIr} constructs a CLR \typename{Exception} with:

\begin{itemize}[leftmargin=2em, noitemsep]
  \item The IR type name and \code{message} field (if present) as the message.
  \item A formatted ``IR Stack'' trace listing each frame's declaring type,
        method name, instruction index, and source location extracted from the
        instruction operand.
\end{itemize}

Source locations are extracted from instruction operands via \code{GetLocationFromOperand},
which probes the \texttt{location}, \texttt{loc}, and \texttt{sourceLocation} JSON
properties and normalises line number fields (\texttt{line}, \texttt{startLine},
\texttt{start}).

% ═══════════════════════════════════════════════════════════════════════════════
\section{Unity Object Wrapping \& Unwrapping}
% ═══════════════════════════════════════════════════════════════════════════════

\begin{longtable}{lp{9cm}}
  \toprule
  \textbf{Helper} & \textbf{Description} \\
  \midrule
  \endhead
  \code{WrapGo(GameObject?)}          & Creates a \typename{ManagedObject} with
                                        type \code{"UnityEngine.GameObject"},
                                        embedding the CLR instance and mirroring
                                        \code{name}, \code{tag},
                                        \code{activeSelf} as IR fields. \\
  \code{UnwrapGo(object?)}            & Extracts the \typename{GameObject} from
                                        either a direct CLR reference or a
                                        \typename{ManagedObject} with
                                        \code{"clrInstance"} metadata. \\
  \code{WrapTransform(Transform?)}    & Creates a \typename{ManagedObject} for a
                                        \code{Transform}, embedding the CLR
                                        instance. \\
  \code{GetTransform(object?)}        & Extracts a \typename{Transform}: from a
                                        \typename{ManagedObject} wrapping a
                                        Transform directly, or from one wrapping a
                                        \typename{GameObject}. \\
  \code{WrapComponent(Component?)}    & Creates a \typename{ManagedObject} using
                                        the component's actual CLR type name. \\
  \code{WrapUnityObject(Object?)}     & Dispatches to \code{WrapGo},
                                        \code{WrapComponent}, or a generic wrap. \\
  \code{UnwrapUnityObject(object?)}   & Extracts any \code{UnityEngine.Object}. \\
  \code{GetClr<T>(object?)}           & Generic helper; returns \code{T} from
                                        direct value or \code{"clrInstance"}
                                        metadata. \\
  \code{SetClrInstance(object?, object)} & Writes a CLR instance into the
                                           \code{"clrInstance"} metadata slot
                                           of a \typename{ManagedObject}. \\
  \bottomrule
\end{longtable}

\subsection{Struct Boxing Helpers}

\begin{longtable}{ll}
  \toprule
  \textbf{Method} & \textbf{Produces} \\
  \midrule
  \code{BoxV3(Vector3)}     & \typename{ManagedObject}(\code{"UnityEngine.Vector3"}) with
                              \code{x,y,z} fields \\
  \code{BoxV2(Vector2)}     & \typename{ManagedObject}(\code{"UnityEngine.Vector2"}) with
                              \code{x,y} fields \\
  \code{BoxQ(Quaternion)}   & \typename{ManagedObject}(\code{"UnityEngine.Quaternion"})
                              with \code{x,y,z,w} fields \\
  \code{BoxColor(Color)}    & \typename{ManagedObject}(\code{"UnityEngine.Color"}) with
                              \code{r,g,b,a} fields \\
  \bottomrule
\end{longtable}

Corresponding unboxing helpers (\code{ToVector3}, \code{ToVector2},
\code{ToQuaternion}, \code{ToColor}) first attempt to read the \code{"clrInstance"}
metadata; if absent they reconstruct the struct from the individual component fields.

% ═══════════════════════════════════════════════════════════════════════════════
\section{Thread Safety}
% ═══════════════════════════════════════════════════════════════════════════════

\begin{itemize}[leftmargin=1.5em]
  \item The \emph{process-global} CLR type registry (\code{s\_registeredClrTypes})
        is populated once under \code{s\_regLock} and then read without locking.
        It is safe to call from any thread after the first instantiation.
  \item A single \typename{UnityIRRuntime} \emph{instance} is \textbf{not}
        thread-safe. Each \typename{UnityIRRuntime} instance owns its own call stack,
        static fields dictionary, and native method table. Do not share an instance
        across threads without external synchronisation.
\end{itemize}

% ═══════════════════════════════════════════════════════════════════════════════
\section{Assembly Definition}
% ═══════════════════════════════════════════════════════════════════════════════

The runtime is shipped as the Unity assembly definition
\texttt{Gamerboy030607.Objectirunity} located at\\
\texttt{Packages/UnityIR/Runtime/Gamerboy030607.Objectirunity.asmdef}.

Dependencies visible to user assemblies:
\begin{itemize}[leftmargin=2em, noitemsep]
  \item \code{ObjectIR.AST} (provides \typename{TextIrParser}, \typename{ModuleNode},
        \typename{ClassNode}, etc.)
  \item \code{ObjectIR.FobCompiler} (provides \typename{FobIrReader})
  \item \code{UnityEngine} (implicit in all Unity projects)
  \item \code{System.Text.Json} (.NET standard library)
\end{itemize}

\end{document}
